\begin{table}[t]
\centering\small
\setlength{\tabcolsep}{4pt}
\caption{Số lần tính hàm trong tối ưu, xuất phát từ điểm đều trên đơn hình.
AC là phương án so sánh toàn lồi; ``lấy mẫu'' và ``tinh chỉnh'' là hai biến thể của SG-dyn.
$\dagger$ nghĩa là không đạt sai số cho phép trong giới hạn số điểm.
Cột cuối là số điểm mà DTAP vi phạm ràng buộc thứ ba, chia cho
số điểm đã tính của nó.}
\label{tab:activation-tuning}
\begin{tabular}{@{}lrrrrrl@{}}
\toprule
Bộ dữ liệu & DTAP & AC & SG-const & Lấy mẫu & Tinh chỉnh & Hoạt động của $F_3$\\
\midrule
Glass        & 89 & 89 & 520 & $1000^\dagger$ & 100 & $0/89$\\
Diabetes     & 14 & 14 & 350 & $1000^\dagger$ & 51 & $0/14$\\
Segmentation & 3 & 3 & 124 & 170 & 3 & $0/3$\\
Abalone      & 48 & 48 & 638 & $1000^\dagger$ & 3 & $1/48$\\
Satimage     & 2 & 2 & 60 & 102 & 2 & $0/2$\\
Yeast        & 4 & 4 & 324 & 305 & 3 & $0/4$\\
Spambase     & 5 & 5 & 93 & 295 & 56 & $0/5$\\
Pendigits    & 2 & 2 & 49 & 10 & 31 & $0/2$\\
\midrule
Movement Libras & 637 & 637 & 62 & $1000^\dagger$ & 144 & $0/637$\\
\bottomrule
\end{tabular}
\end{table}

Mỗi khi $F_3(w^k)\le b_3$ dọc theo một lần chạy, tập active của nhánh giả lồi rỗng,
$\mu_k=1$ và DTAP trùng từng bước với phương án so sánh toàn lồi;
trên tám bài toán có mức hoạt động của $F_3$ bằng không, hai cột DTAP và AC giống hệt nhau
phản ánh sự trùng khớp này. Trên Abalone, $F_3$ chỉ bị vi phạm tại điểm xuất phát
đều ($1/48$), nơi $J^{\ps}_0\approx1.56\times10^{-3}$ rất nhỏ so với
$J^{\cvx}_0\approx366.3$, nên $1-\mu_0\approx4.27\times10^{-6}$; các giá trị của $J$
dọc hai quỹ đạo thỏa $|J_{\rm DTAP}-J_{\rm AC}|\le3.18\times10^{-5}\,J_{\rm AC}$ tại
mọi điểm đã tính,
và cả hai lần chạy dừng sau $48$ điểm đã tính. Vì vậy các cột này
không phải là phép so sánh hai cách gán; hai cách gán chỉ khác nhau về số lần thành công
trên các điểm xuất phát ngẫu nhiên ($26/36$ so với $36/36$).

Việc tinh chỉnh nâng số lần thành công của SG-dyn, xuất phát từ điểm chuẩn, từ $5/9$ lên $9/9$.
SG-dyn tinh chỉnh có thể tốt hơn DTAP, chẳng hạn trên Abalone, và SG-const nhanh hơn
trên Movement Libras. Trên cả bốn điểm xuất phát của cả chín bài toán đặc trưng,
số lần thành công là $26/36$ cho DTAP, $36/36$ cho phương án so sánh toàn lồi,
$33/36$ cho SG-const, $15/36$ cho SG-dyn lấy mẫu, và $28/36$
cho SG-dyn tinh chỉnh. Một biến thể DTAP cân bằng cũng đạt $26/36$; nó gán
trọng số bằng nhau khi cả hai nhánh bị vi phạm, và ngược lại gán toàn bộ trọng số cho nhánh duy nhất bị vi phạm.
Như vậy các phép thử này không chứng tỏ được ưu thế đồng đều của trọng số tỉ lệ,
cũng không chứng tỏ được tính ổn định đối với điểm khởi đầu.

Chín lần tìm kiếm tham số trên các bài toán đặc trưng cần tổng cộng $179\,744$ lần tính hàm trong tối ưu
và $10.80$ giây thời gian giải cho việc tinh chỉnh. Các chi phí này
là phần thêm vào so với Bảng~\ref{tab:activation-tuning}; DTAP không được tinh chỉnh.
Việc lấy mẫu mô-đun ban đầu, tiền xử lý và chi phí phụ của việc tìm kiếm không được tính
vào tổng thời gian giải này. Hồ sơ bổ sung ghi chi phí của
từng lần tìm kiếm.

\paragraph*{Một ràng buộc giả lồi thực sự không lồi.}
Để kiểm tra trực tiếp nhánh giả lồi, xét $K=[-2,2]^d$ và
\begin{equation}\label{eq:logball-test}
\begin{aligned}
 F_1(x)&=\|x+e_1\|^2\le1,\qquad
 F_2(x)=\|x-e_2\|^2\le\tfrac94,\\
 F_3(x)&=\kappa\log(1+\|x-e_1\|^2)\le \kappa\log3.
\end{aligned}
\end{equation}
Ở đây $d\in\{2,10,50\}$ và $\kappa\in\{0.3,1,3\}$ cho chín cấu hình.
Hàm thứ ba là giả lồi: nó là hợp của một hàm bậc hai lồi với một hàm khả vi tăng
có đạo hàm dương. Nó không
lồi, vì giá trị riêng theo phương bán kính của ma trận Hesse là
$2\kappa(1-r^2)/(1+r^2)^2<0$ với $r=\|x-e_1\|>1$.
Điểm minh chứng chính xác $x=-e_1/10$ có các độ chênh dương
$19/100$, $31/25$ và $\kappa\log(300/221)$.
Điểm xuất phát chuẩn là $1.5\mathbf1$; ba điểm xuất phát dùng để đánh giá và ba
điểm xuất phát dùng để tinh chỉnh được lấy đều trên $[0.8,1.8]^d$ với các seed nêu trên
(NumPy \texttt{default\_rng}).
Để lấy mẫu mô-đun, ta dùng $500$ điểm đều trong $K$ (seed $42$).
SG-const dùng cùng bước $\nu_k=10^{-3}$ như trên các bài toán đặc trưng. SG-const
tinh chỉnh dùng một bước hằng $\nu$ chọn từ lưới
$\{10^{-3},10^{-2.5},\ldots,10^1\}$ với các điểm xuất phát tinh chỉnh và quy tắc chọn của
SG-dyn tinh chỉnh; nó chọn $\nu=10^{-0.5}$ ở cả chín cấu hình. Biến thể cân bằng
là biến thể đã mô tả ở trên.

\begin{table}[t]
\centering\small
\setlength{\tabcolsep}{5pt}
\caption{Các phép thử giả lồi không lồi, nhóm theo $\kappa$.
Mỗi ô là số lần thành công trên $12$ lần chạy, kèm trung vị số
điểm đã tính trong ngoặc, kể cả các lần chạy hết số điểm tối đa. Mỗi nhóm
gồm cả ba số chiều và bốn điểm xuất phát. SG-const tinh chỉnh dùng bước
được chọn như mô tả trong bài.}
\label{tab:synthetic-scale}
\begin{tabular}{@{}lrrrrrr@{}}
\toprule
$\kappa$ & DTAP & Cân bằng & SG-const & SG-const tinh chỉnh & SG-dyn lấy mẫu & SG-dyn tinh chỉnh\\
\midrule
$0.3$ & $0\ (1000)$ & $0\ (1000)$ & $0\ (1000)$ & $12\ (15)$ & $8\ (725.5)$ & $12\ (142.5)$\\
$1$   & $12\ (70)$ & $12\ (395.5)$ & $0\ (1000)$ & $12\ (15)$ & $12\ (301.5)$ & $12\ (44)$\\
$3$   & $12\ (4)$ & $12\ (4)$ & $0\ (1000)$ & $12\ (16)$ & $12\ (167)$ & $12\ (11)$\\
\bottomrule
\end{tabular}
\end{table}

DTAP gặp vi phạm ràng buộc thứ ba ở mọi lần trong $36$
lần chạy này. Tính gộp lại, ràng buộc thứ ba bị vi phạm tại $12\,835/12\,877$ điểm đã tính.
Tuy vậy, nó chỉ đạt sai số cho phép ở $24/36$ lần chạy, so
với $32/36$ của SG-dyn lấy mẫu, $36/36$ của SG-dyn tinh chỉnh và $36/36$ của SG-const
tinh chỉnh (Bảng~\ref{tab:synthetic-scale}).
SG-const tinh chỉnh đạt sai số cho phép ở mọi lần chạy với trung vị $15$--$16$
điểm đã tính, ít hơn mọi phương pháp khác tại $\kappa=0.3$ và $\kappa=1$; tại
$\kappa=3$, DTAP và biến thể cân bằng cần trung vị $4$. Các số đếm này
không gồm $147\,331$ lần tính hàm trong việc tìm kiếm tham số của nó.

Với $\nu_k=10^{-3}$, SG-const di chuyển nhiều nhất $999\times10^{-3}=0.999$ so với điểm xuất phát
trong giới hạn số điểm, trong khi mọi điểm xuất phát cách
$\Omega$ ít nhất $1.11$; vì vậy các số đếm bằng không của nó do cỡ bước và giới hạn số điểm quyết định.
Với $\kappa\le1$, nhánh lồi được thỏa từ điểm đã tính thứ hai trở đi,
nên sau đó DTAP đi theo phép lặp một nhánh của Hệ quả~\ref{cor:ps}; cả hai nhánh
cùng bị vi phạm chỉ tại $32$ trong $12\,877$ điểm đã tính này.

Mọi thất bại của DTAP đều xảy ra tại $\kappa=0.3$. Thay đổi $\kappa$ không làm thay đổi tập chấp nhận được,
nhưng $v_3$ và $\nabla F_3$ đều tỉ lệ với $\kappa$, nên khi
$\|g^{\ps}_k\|\le1$ bước giả lồi $(1-\mu_k)\gamma_k g^{\ps}_k$ tỉ lệ với
$\gamma_k\kappa^2$. Bảng~\ref{tab:dtap-sweep} lần lượt thay đổi từng tham số của DTAP
với $\kappa\in\{0.3,1\}$. Với $\gamma_k=a(k+0.5)^{-0.55}$, $a=10$ và $a=100$ cho
$12/12$ lần thành công tại $\kappa=0.3$; $\rho_{\cvx}=1$ cho $4/12$, đều với $d=2$;
$\delta=10^{-1}$ tái tạo chính xác các điểm lặp của cấu hình gốc (Nhận xét~\ref{rem:delta-role}); các
bước giảm nhanh hơn $s=0.75$ và $s=1$ không thành công lần nào tại $\kappa=0.3$, và
$s=1$ cũng thất bại cả $12$ lần chạy tại $\kappa=1$. Vì vậy các thất bại tại $\kappa=0.3$
phù hợp với việc bước giảm dần quá nhỏ so với thang của ràng buộc
này, chứ không phải với một thay đổi về tính chấp nhận được. Các phép thử này thay đổi
từng tham số một trên một họ bài toán và không chỉ ra được một cấu hình
tốt nhất nói chung.
Trên~\eqref{eq:logball-test}, $\delta=10^{-1}$ và $\delta=10^{-3}$ cho các dãy lặp giống hệt nhau vì
nhánh giả lồi chỉ có một chỉ số, tại mọi điểm xuất phát hoặc chỉ $F_1$ bị vi phạm hoặc
$v_1-v_2\ge4.79$, và nhánh lồi được thỏa từ điểm đã tính thứ hai trở đi; vì vậy
dòng $\delta$ của Bảng~\ref{tab:dtap-sweep} là một phép kiểm tra tính nhất quán, không phải bằng chứng về
vai trò của $\delta$.
Các lần tìm kiếm tham số trên bài toán tổng hợp tốn $188\,295$ lần tính hàm trong tối ưu và
$10.98$ giây thời gian giải cho SG-dyn, và $147\,331$ lần tính hàm và $8.27$
giây cho SG-const, với cùng các khoản loại trừ như trên, trong môi trường
nêu dưới đây.

\begin{table}[t]
\centering\small
\setlength{\tabcolsep}{5pt}
\caption{DTAP trên~\eqref{eq:logball-test} khi thay đổi từng tham số một,
$\gamma_k=a(k+0.5)^{-s}$. Cấu hình gốc $\rho_{\cvx}=1.98$, $\delta=10^{-3}$, $a=1$,
$s=0.55$ là cột DTAP của Bảng~\ref{tab:synthetic-scale}. Mỗi ô là
số lần thành công trên $12$ lần chạy (ba số chiều, bốn điểm xuất phát), kèm
trung vị số điểm đã tính trong ngoặc, kể cả các lần chạy hết số điểm tối đa.}
\label{tab:dtap-sweep}
\begin{tabular}{@{}lrr@{}}
\toprule
Cấu hình & $\kappa=0.3$ & $\kappa=1$\\
\midrule
Cấu hình gốc & $0\ (1000)$ & $12\ (70)$\\
$\rho_{\cvx}=1$ & $4\ (1000)$ & $12\ (223.5)$\\
$\delta=10^{-1}$ & $0\ (1000)$ & $12\ (70)$\\
$a=10$ & $12\ (77)$ & $12\ (6)$\\
$a=100$ & $12\ (6)$ & $12\ (48)$\\
$s=0.75$ & $0\ (1000)$ & $12\ (374)$\\
$s=1$ & $0\ (1000)$ & $0\ (1000)$\\
\bottomrule
\end{tabular}
\end{table}

\paragraph*{Khả năng tái lập.}
Kết quả chọn đặc trưng được lưu với Python~3.11.7 và NumPy~1.26.4
trên Windows, với yêu cầu một luồng BLAS/OpenMP. Các hàm tính giá trị (oracle) cho bài toán đặc trưng dùng dạng toàn phương
nén theo đại số; sai khác đã kiểm tra so với giá trị và gradient
gốc lần lượt không quá $1.82\times10^{-12}$ và
$1.10\times10^{-11}$. Thời gian giải báo cáo cho các bài toán đặc trưng ứng
với cách cài đặt nén này. Gói mã đi kèm chỉ chứa các hàm tính giá trị không nén;
nó tái tạo các số lần tính hàm, còn thời gian giải được báo cáo là của
cách cài đặt nén ban đầu. Các Bảng~\ref{tab:synthetic-scale}
và~\ref{tab:dtap-sweep} và chi phí tinh chỉnh trên bài toán tổng hợp được tạo với
Python~3.13.4 và NumPy~2.1.3 trên Windows với số luồng BLAS
mặc định; các số đếm của chúng là tất định, còn thời gian giải là các
phép đo một lần trên máy đó. Với Python~3.13.4 và NumPy~2.1.3, gói mã còn
tái tạo chính xác mọi giá trị số trong các Bảng~\ref{tab:datasets}--\ref{tab:wilcoxon}
(Bảng~\ref{tab:perf-summary} và~\ref{tab:wilcoxon} từ điểm chưa làm tròn) và trong
Bảng~\ref{tab:activation-tuning}, cùng các số đếm nêu ở Mục~\ref{ssec:results}
và~\ref{ssec:feasibility-audit}.
Các script bổ sung, quỹ đạo thô, hồ sơ tinh chỉnh, điểm minh chứng,
mã băm của mã nguồn và hướng dẫn kiểm tra được cung cấp kèm theo kết quả.

\section{Kết luận}\label{sec:conclusion}
DTAP giải các bài toán chấp nhận tách phi tuyến với hộp chuẩn và các thành phần lồi và giả lồi, bằng cách kết hợp một bước lồi điều chỉnh theo phần dư với một bước giả lồi có bảo vệ. Dưới giả thiết nhất quán và khả vi liên tục, toàn bộ dãy hội tụ tới một điểm chấp nhận được; không có mô-đun chính quy nào là đầu vào, mặc dù phân tích dùng các cận trên một tập compact chứa dãy lặp. Trong ứng dụng chọn đặc trưng, mọi ràng buộc đều lồi, và các bài toán thử có ràng buộc giả lồi không lồi kiểm tra nhánh giả lồi. Trong thực nghiệm, các phương pháp dưới gradient đã tinh chỉnh có thể bằng hoặc tốt hơn DTAP, và hiệu năng của DTAP với số vòng lặp hữu hạn phụ thuộc vào điểm xuất phát và thang cỡ bước. Tốc độ cho toàn bộ phần dư $J(x^k)$, ngoài cận của nhánh lồi ở Nhận xét~\ref{rem:rate}, và các bảo đảm dưới điều kiện cận sai số vẫn còn là vấn đề mở.

\section*{Tuyên bố về xung đột lợi ích}
Các tác giả tuyên bố không có xung đột lợi ích tài chính hay quan hệ cá nhân nào đã biết có thể ảnh hưởng đến công trình được báo cáo trong bài.
\section*{Tài trợ}
Các tác giả không nhận tài trợ riêng nào cho công trình này.
\section*{Tuyên bố về việc dùng AI tạo sinh}
OpenAI Codex đã được dùng trong quá trình chỉnh sửa để biên tập ngôn ngữ.
\section*{Tính sẵn có của dữ liệu}
Các bộ dữ liệu có tại UCI Machine Learning Repository~\citep{kelly2023uci}. Gói tái lập đi kèm chứa các script, kết quả số đã ghi lại, và hướng dẫn tạo lại các bảng đã báo cáo.

\section*{ORCID}
\noindent Nguyen Van The~\raisebox{-0.1ex}{\textit{https://orcid.org/0009-0004-6548-420X}}\\
Tran Ngoc Thang~\raisebox{-0.1ex}{\textit{https://orcid.org/0009-0000-9549-5736}}\\
Nguyen Hong Son~\raisebox{-0.1ex}{\textit{https://orcid.org/0009-0006-3106-8296}}


